\section{Modules}\label{sec:modules}

\paragraph{Modules over rings}

\begin{definition}\label{def:endomorphism_ring}\mcite[102]{Bourbaki1998Algebra1to3}
  We define the \term[ru=кольцо эндоморфизмов (\cite[130]{Курош1973ОбщаяАлгебра})]{endomorphism ring} of an \hyperref[def:abelian_group]{abelian group} \( G \) as its \hyperref[def:endomorphism_semiring]{endomorphism semiring} \( \End(G) \) endowed with the obvious additive inverses.
\end{definition}

\begin{definition}\label{def:module}
  A left (resp. right) \( R \)-\term[bg=модул (\cite[27]{КоцевСидеров2016КомутативнаАлгебра}, ru=модуль (\cite[def. 9.3.1]{Винберг2014КурсАлгебры})]{module} is a left (resp. right) \( R \)-\hyperref[def:semimodule]{semimodule} over a \hyperref[def:ring]{ring} endowed with the obvious additive inverses.

  We give formal definitions based on those for semimodules:
  \begin{thmenum}[series=def:module]
    \thmitem{def:module/action}\mcite[156]{Aluffi2009Algebra} A left \( R \)-module is a \hyperref[def:ring/morphism]{ring homomorphism} from \( R \) to the \hyperref[def:endomorphism_ring]{endomorphism ring} \( \End(M) \). A right \( R \)-module is a homomorphism from \( R \) to the \hyperref[def:semiring/opposite]{opposite ring} \( \End(M)^\oppos \).

    \thmitem{def:module/operation}\mcite[def. III.5.2]{Aluffi2009Algebra} Similarly to \( R \)-semimodules, a left \( R \)-module can be described via a scalar multiplication operation \( \cdot: R \times M \to M \) satisfying the conditions from \cref{def:semimodule/operation}.

    \Cref{thm:module_absorption_condition} implies that the absorption law \eqref{eq:def:semimodule/operation/absorption} becomes redundant.
  \end{thmenum}

  Modules have the following metamathematical properties:
  \begin{thmenum}
    \thmitem{def:module/theory}\mcite[\S 3.3.3]{TentZiegler2012ModelTheory} Similarly to how we described an \hyperref[def:fol_equational_theory]{equational theory} for left \( R \)-semimodules in \cref{def:semimodule/theory}, for \( R \)-modules we can also use a family \( \seq{ \synscalar_r }_{r \in R} \) of unary function symbols.

    We can extend the theory of \( R \)-semimodules with an additional binary subtraction symbol \( \synminus \), for which we add the left and right inverse axioms \eqref{eq:def:additive_semigroup/theory/inverse/left} and \eqref{eq:def:additive_semigroup/theory/inverse/right}.

    \Cref{thm:module_absorption_condition} implies that the absorption formulas \eqref{eq:def:semimodule/theory/absorption} becomes unnecessary as an axiom because it can be derived from the rest.

    Alternatively, we can \hyperref[def:many_sorted_fol]{two-sorted first-order} theory described in \cref{rem:two_sorted_semimodules}.

    \thmitem{def:module/morphism}\mcite[160]{Aluffi2009Algebra} We say that the function \( \varphi: M \to N \) between (left or right) \( R \)-modules is a \term{homomorphisms of \( R \)-modules} if it it if preserves vector addition and scalar multiplication with all elements of \( R \). These coincide with the \hyperref[def:linear_function]{\( R \)-linear functions}.

    The conditions stated are equivalent to those given in \cref{def:semimodule/morphism} for \( R \)-semimodule homomorphisms, however the reasoning that these homomorphisms coincide with the corresponding \hyperref[def:fol_homomorphism]{first-order homomorphism} is different --- see \cref{thm:module_homomorphism_is_fol_homomorphism}.

    \thmitem{def:module/category}\mcite[158]{Aluffi2009Algebra} For a fixed \hyperref[def:grothendieck_universe]{Grothendieck universe} \( \mscrU \) and a fixed ring \( R \), the first-order theories of \( R \)-modules give rise to several \hyperref[def:fol_theory_model_functor/obj]{categories of \( \mscrU \)-small models}.

    We denote the category of left (resp. right) \( R \)-modules by \( \ucat[R]{Mod} \) (resp. \( \ucat{Mod}[R] \)). Variations of this notation are discussed in \cref{rem:category_of_modules_notation}.

    It is a very well-behaved category, even more than the category \hyperref[def:group/category]{\( \ucat{Grp} \)} of \( \mscrU \)-small groups.
    \begin{itemize}
      \item Just as with \( R \)-semimodules, \fullref{thm:free_semimodule_universal_property} implies that \hyperref[def:free_semimodule]{free modules} provide a \hyperref[def:category_adjunction]{left adjoint} to the \hyperref[def:fol_theory_forgetful_functor/set]{forgetful functor}
      \begin{equation*}
        U: \ucat[R]{Mod} \to \ucat{Set}.
      \end{equation*}

      \item With the above left adjoint, \cref{thm:concrete_category_function_invertibility/injective} implies that an \( R \)-module homomorphism is a \hyperref[def:morphism_invertibility/left_cancellative]{monomorphism} if and only if it is \hyperref[def:function_invertibility/injective]{injective}.

      \item Every epimorphism in \( \ucat[R]{Mod} \) is surjective. This will be shown in \cref{thm:module_epimorphisms_are_surjective}.

      \item The zero \( R \)-semimodules, which we will define in \cref{def:semimodule/trivial}, are \hyperref[def:universal_objects/zero]{zero objects} in the categories \( \ucat[R]{Mod} \) and \( \ucat{Mod}[R] \).

      \item \hyperref[def:categorical_quotient_object]{Quotient objects} can be described via plain \hyperref[def:module/subobject]{submodules} (without additional like normal subgroups or ring ideals).
    \end{itemize}

    \thmitem{def:module/subobject}\mcite[160]{Aluffi2009Algebra} An \( R \)-\term[bg=подмодул (\cite[239]{ГеновМиховскиМоллов1991Алгебра}), ru=подмодуль (\cite[396]{Винберг2014КурсАлгебры})]{submodule} of \( M \) is an \hyperref[def:group/subobject]{subgroup} that is closed under scalar multiplication with all elements of \( R \).

    An \( R \)-submodule can also be characterized as an \( R \)-\hyperref[def:semimodule/subobject]{sub-semimodule} that is closed under additive inverses. In practice, it is more easily characterized as a subset closed under \hyperref[def:linear_combination]{linear combinations} with coefficients from \( R \).

    It follows from \cref{thm:fol_equational_theory_models/substructure} that the submodules of \( M \) are precisely its \hyperref[def:fol_substructure/submodel]{first-order submodels}, and from \cref{thm:fol_categorical_subobjects} that they are its \hyperref[def:categorical_subobject]{categorical subobjects}.

    Furthermore, as a consequence of \cref{thm:fol_equational_theory_models/image}, the image of a linear map is a \( R \)-submodule of its codomain.

    Submodules are well-behaved because they correspond \emph{exactly} to first-order congruences --- see \cref{thm:submodules_and_congruences}.

    \thmitem{def:module/generated}\mcite[204]{Bourbaki1998Algebra1to3} We define the \term{generated \( R \)-submodule} \( \lin_R A \) of the subset \( A \) of the \( R \)-module \( M \) as the as the intersection of \( R \)-submodules of \( M \) containing \( A \). It coincides with the first-order submodel of \( M \) \hyperref[def:fol_generated_substructure]{generated} by \( A \).

    As shown in \cref{thm:generated_submodule}, the generated sub-semimodule can be characterized both \hyperref[def:linear_combination]{linear combinations} and \hyperref[def:linear_span]{linear spans}.

    \thmitem{def:module/trivial}\mcite[196]{Bourbaki1998Algebra1to3} We call the \hyperref[def:group/trivial]{trivial additive group} \( \set{ 0 } \) the \term[en=trivial module (\cite[176]{Aluffi2009Algebra})]{trivial \( R \)-module} or \term{zero \( R \)-module}.

    In accordance with \cref{def:trivial_object}, due to \cref{thm:equational_theory_zero_morphisms/trivial}, it is a \hyperref[def:universal_objects/zero]{zero object} in the categories \( \ucat[R]{Mod} \) and \( \ucat{Mod}[R] \).

    \thmitem{def:module/kernel}\mcite[197]{Bourbaki1998Algebra1to3} We define the \term[bg=ядро (\cite[241]{ГеновМиховскиМоллов1991Алгебра}), ru=ядро (\cite[397]{Винберг2014КурсАлгебры})]{kernel} of an \( R \)-module homomorphism \( \varphi: M \to N \) as the set
    \begin{equation}\label{eq:def:module/kernel}
      \ker \varphi \coloneqq \varphi^{-1}(0_N) = \set{ x \in M \given \varphi(x) = 0_N }.
    \end{equation}

    \Cref{thm:equational_theory_zero_morphisms/kernel} implies that this is a submodule acting as the \hyperref[def:categorical_kernels/kernel]{categorical kernel} of \( \varphi \).

    In vector spaces, kernels are also called \enquote{null spaces} --- see \cref{def:vector_space/kernel}.

    \thmitem{def:module/quotient}\mimprovised \hyperref[def:fol_congruence]{First-order congruences} in the theory of modules are even better behaved than those for groups and rings. \Cref{thm:submodules_and_congruences} shows that congruences correspond \emph{exactly} to submodules.

    For this reason, a \term[bg=фактор-модул (\cite[241]{ГеновМиховскиМоллов1991Алгебра}), ru=фактормодуль (\cite[398]{Тыртышников2017ОсновыАлгебры}), en=quotient module (\cite[197]{Bourbaki1998Algebra1to3})]{quotient group} \( M / N \) is more conveniently defined with respect to a submodule \( N \) rather than with respect to the corresponding congruence.

    \thmitem{def:module/bimodule}\mcite[def. 3.5]{Jacobson1989BasicAlgebraII} Similarly to \hyperref[def:semimodule/bisemimodule]{bisemimodules}, we define an \( (R, S) \)-\term{bimodule} as a triple \( (R, S, M) \), where \( M \) is an abelian group that is simultaneously a left \( R \)-module and a right \( S \)-module, and the associativity-like condition \eqref{eq:def:semimodule/bisemimodule/associativity} holds.
  \end{thmenum}
\end{definition}

\begin{remark}\label{rem:category_of_modules_notation}
  As mentioned in \cref{def:module/category}, there are several established notations for left and right \( R \)-modules.

  \begin{itemize}
    \item To distinguish the categories of left and right \( R \)-modules, it is convenient to write \( \cat{R\mathdash Mod} \) and \( \cat{Mod\mathdash R} \).

    This is the more popular convention, introduced in
    \cite[12]{Jacobson1985BasicAlgebraI},
    \cite[158]{Aluffi2009Algebra} and
    \cite[12]{MacLane1998CategoryTheory}.

    \Textcite[446]{Rotman2015AdvancedModernAlgebraPart1} denotes these categories by \( \cat{{}_R Mod} \) and \( \cat{Mod_R} \).

    \item Another convention is to add the ring as a suffix and write \( \cat{Mod}_R \). It is less popular because it has the downside of not distinguishing between left and right \( R \)-modules.

    Thus, the convention is less popular. It is introduced in \cite[example 1.1.3]{Riehl2016CategoryTheory} and used through the book.

    For vector spaces, where no distinction between left and right scalar multiplication is necessary, the notation \( \cat{Vect}_\BbbK \) is used in
    \cite[xiii]{Leinster2014BasicCategoryTheory}
  \end{itemize}
\end{remark}

\begin{proposition}\label{thm:module_absorption_condition}
  The \hyperref[def:semimodule]{semimodule} absorption condition \eqref{eq:def:semimodule/operation/absorption} is redundant in \hyperref[def:module]{modules} (over unital rings).
\end{proposition}
\begin{proof}
  \begin{equation*}
    0_R \cdot x
    =
    (-1_R + 1_R) \cdot x
    \reloset {\eqref{eq:def:semimodule/operation/scalar_addition_distributivity}} =
    -1_R \cdot x + 1_R \cdot x
    \reloset {\eqref{eq:def:semimodule/operation/action/identity}} =
    -x + x
    =
    0_M.
  \end{equation*}
\end{proof}

\begin{proposition}\label{thm:module_homomorphism_is_fol_homomorphism}
  The \hyperref[def:module/morphism]{module homomorphisms} preserve zeros and additive inverses, i.e. are \hyperref[def:fol_homomorphism]{first-order homomorphisms}.
\end{proposition}
\begin{proof}
  Follows from \cref{thm:group_homomorphism_single_condition}.
\end{proof}

\begin{proposition}\label{thm:abelian_group_is_module}
  The \hyperref[def:category]{categories} \hyperref[def:abelian_group]{of abelian groups} and \hyperref[def:module/category]{of \( \BbbZ \)-semimodules} are \hyperref[rem:category_similarity/isomorphism]{isomorphic}.

  More concretely, every abelian group \( M \) is a left \( \BbbZ \)-module with scalar multiplication given by the recursively defined multiplication operation \eqref{eq:def:additive_semigroup/iterated}.

  Conversely, in every left \( \BbbZ \)-module, scalar multiplication matches the recursively defined multiplication.
\end{proposition}
\begin{comments}
  \item This result specializes a more general statement for semimodules over semirings in \cref{thm:commutative_monoid_is_semimodule} and specializes to the case for algebras and integers in \cref{thm:ring_is_integer_algebra}.
\end{comments}
\begin{proof}
  Simple refinement of \cref{thm:commutative_monoid_is_semimodule}.
\end{proof}

\paragraph{Module quotients}

\begin{proposition}\label{thm:submodules_and_congruences}
  Fix a (left or right) \( R \)-\hyperref[def:module]{module} \( M \). \hyperref[def:fol_congruence]{First order congruences} on \( M \) and \hyperref[def:module/subobject]{\( R \)-submodules} of \( M \) are related as follows:
  \begin{thmenum}
    \thmitem{thm:submodules_and_congruences/cong_to_submodule} For every congruence in \( M \), the coset \( [0_M] \) of the zero element is an \( R \)-submodule of \( M \).

    \thmitem{thm:submodules_and_congruences/submodule_to_cong} Conversely, for every \( R \)-submodule \( N \), the relation
    \begin{equation}\label{eq:thm:submodules_and_congruences/submodule_to_cong}
      x \cong y \T{if and only if} x - y \in N,
    \end{equation}
    is a congruence.

    \thmitem{thm:submodules_and_congruences/inverse} The procedures for obtaining a congruence from a submodule and vice versa are inverses.
  \end{thmenum}
\end{proposition}
\begin{comments}
  \item This allows us to consider the quotient \( M / N \) with respect to an \( R \)-submodule rather than \( M / {\cong} \) with respect to a congruence.

  \item This is based on a similar statement for \hyperref[def:group]{groups} --- see \cref{thm:normal_subgroups_and_congruences}. Another similar statement holds for \hyperref[def:ring]{rings} --- see \cref{thm:ring_ideals_and_congruences} --- and for \hyperref[def:algebra_over_ring]{algebras over rings} --- see \cref{thm:algebra_ideals_and_congruences} and for \hyperref[def:algebra/ideal]{algebra ideals} --- see \fullref{thm:lattice_theorem_for_algebra_ideals}.
\end{comments}
\begin{proof}
  \SubProofOf{thm:submodules_and_congruences/cong_to_submodule} Since module congruences are also congruences of the additive group, \cref{thm:normal_subgroups_and_congruences/cong_to_subgroup} implies that \( [0_M] \) is a \hyperref[def:normal_subgroup]{normal subgroup} of the additive group of \( M \).

  We must show that \( [0_M] \) is closed under scalar multiplication. This follows from \Cref{thm:semimodule_vector_absorption}, which states implies that \( 0_M \) absorbs scalars, i.e. \( r \cdot 0_M = 0_M \) for every \( r \) in \( R \).

  \SubProofOf{thm:submodules_and_congruences/submodule_to_cong} Similarly, since every \( R \)-submodule is a normal subgroup of the additive group, it follows from \cref{thm:normal_subgroups_and_congruences/subgroup_to_cong} that \( {\cong} \) defined as \( x \cong y \) if \( x - y \in I \) is a congruence on the additive group.

  To prove that it is an \( R \)-module congruence, we must show that \( x \cong y \) implies \( rx \cong ry \) every scalar \( r \). This holds because of distributivity:
  \begin{equation*}
    rx - ry
    \reloset {\eqref{eq:def:semimodule/operation/vector_addition_distributivity}} =
    r(\underbrace{x - y}_{N}).
  \end{equation*}

  Since \( N \) is closed under scalar multiplication, we conclude that \( rx - ry \) belongs to \( N \) and hence \( rx \cong ry \).

  \SubProofOf{thm:submodules_and_congruences/inverse} Special case of \cref{thm:normal_subgroups_and_congruences/inverse}.
\end{proof}

\begin{theorem}[Lattice theorem for submodules]\label{thm:lattice_theorem_for_submodules}
  Given an \( R \)-\hyperref[def:module/subobject]{submodule} \( N \) of the \( R \)-module \( M \), the function \( K \mapsto K / N \) is an \hyperref[def:lattice/homomorphism]{isomorphism of complete lattices} between the \hyperref[def:lattice_ideal/principal]{principal filter} of \( R \)-submodules of \( M \) containing \( N \) and the \hyperref[thm:semiring_of_ideals/lattice]{lattice of \( R \)-submodules} of the \hyperref[def:module/quotient]{quotient} \( M / N \).
\end{theorem}
\begin{comments}
  \item See \fullref{thm:lattice_theorem_for_substructures} for a more verbose formulation of this theorem in a general setting.

  \item This is based on a similar statement for \hyperref[def:group/subobject]{subgroups} --- see \fullref{thm:lattice_theorem_for_subgroups}. Another similar statement holds for \hyperref[ring/ideal]{ring ideals} --- see \fullref{thm:lattice_theorem_for_ring_ideals}.
\end{comments}
\begin{proof}
  Fix an \( R \)-submodule \( N \) of \( M \) and let \( x \cong y \) if \( x - y \) is in \( N \).

  Note that every \( R \)-submodule of \( M \) is a subgroup of (the additive group of) \( M \) and that \( {\cong} \) is also a congruence on the additive group of \( M \). \Fullref{thm:lattice_theorem_for_subgroups} allows us to conclude that the \( R \)-submodule \( K \) of \( N \) contains \( N \) if and only if it is compatible with \( {\cong} \).

  The lattice isomorphism then follows from \fullref{thm:lattice_theorem_for_substructures}.
\end{proof}

\begin{proposition}\label{thm:kernel_is_submodule}
  The \hyperref[def:module/kernel]{kernel} of an \( R \)-\hyperref[def:module/morphism]{module homomorphism} is an \( R \)-\hyperref[def:module/subobject]{submodule} of its domain.
\end{proposition}
\begin{proof}
  Follows from \cref{thm:kernel_is_normal_subgroup} and absorption.
\end{proof}

\begin{theorem}[Homomorphism theorem for modules]\label{thm:homomorphism_theorem_for_modules}
  Every \( R \)-\hyperref[def:module/morphism]{module homomorphism} \( {\varphi: M \to N} \) induces an isomorphism
  \begin{equation*}
    \begin{aligned}
      &\Phi: M / \ker \varphi \to \img \varphi, \\
      &\Phi([x]) \coloneqq \varphi(x)
    \end{aligned}
  \end{equation*}
  between the \hyperref[def:module/quotient]{quotient module} \( M / \ker \varphi \) and the \hyperref[def:set_valued_map/image]{set-theoretic image} \( \img \varphi = \varphi[M] \) of \( \varphi \).
\end{theorem}
\begin{comments}
  \item This is a restatement of \fullref{thm:fol_homomorphism_factorization} that uses submodules rather than congruences.

  It is based on a similar statement for \hyperref[def:group]{groups} --- see \fullref{thm:homomorphism_theorem_for_rings}. Another similar statement holds for \hyperref[def:ring]{rings} --- see \fullref{thm:homomorphism_theorem_for_rings} --- and for \hyperref[def:algebra_over_ring]{algebras over rings} --- see \fullref{thm:homomorphism_theorem_for_algebras}.

  \item In the categories of modules, the quotient \( M / \ker \varphi \) is sometimes referred to as a coimage --- see \cref{thm:module_categorical_kernels/coimage}.
\end{comments}
\begin{proof}
  First, note that the domain and codomain of \( \varphi \) are indeed modules --- the quotient \( M / \ker \varphi \) is an \( R \)-module by definition, while the image \( \img \varphi \) is an \( R \)-module as a consequence of \cref{thm:fol_equational_theory_models/image}.

  \Fullref{thm:homomorphism_theorem_for_groups} implies that \( \Phi \) is a well-defined isomorphism of additive groups. In order for \( \Phi \) to be an \( R \)-module isomorphism, we must show that it is well-defined with respect to scalar multiplication, and that it preserves scalar multiplication.

  \Cref{thm:submodules_and_congruences/cong_to_submodule} implies that the relation
  \begin{equation*}
    x \cong y \T{if and only if} x - y \in \ker \varphi
  \end{equation*}
  is a congruence of \( R \)-modules rings (since \( \ker \varphi \) is an \( R \)-submodule by \cref{thm:kernel_is_submodule}). This is the same congruence used in \fullref{thm:homomorphism_theorem_for_groups} to show that \( \Phi \) is an isomorphism of additive groups. In particular, the scalar product \( r[x] \) is \emph{defined} in \( R / \ker \varphi \) as the coset \( [rx] \), so \( \Phi \) preserves scalar multiplication.

  Therefore, \( \Phi \) is a homomorphism of \( R \)-modules. Since it is bijective, it is in fact an isomorphism.
\end{proof}

\paragraph{Cokernels}

\begin{definition}\label{def:module_cokernel}\mcite[96]{Jacobson1989BasicAlgebraII}
  We define the \term{cokernel} of the \( R \)-\hyperref[def:module/morphism]{module homomorphism} \( \varphi: M \to N \) as the \hyperref[def:module/quotient]{quotient}
  \begin{equation*}
    \co\ker \varphi \coloneqq N / \img f.
  \end{equation*}
\end{definition}

\begin{proposition}\label{thm:module_categorical_kernels}
  In the category \hyperref[def:module/category]{\( \ucat[R]{Mod} \)} of left \( R \)-modules or \hyperref[def:ring/category]{\( \ucat{Rng} \)}), \hyperref[def:module/kernel]{kernels} and \hyperref[def:module_cokernel]{cokernels} related to the corresponding categorical notions from \cref{def:categorical_kernels} as follows:
  \begin{thmenum}
    \thmitem{thm:module_categorical_kernels/kernel} The module-theoretic kernel \( \ker \varphi = \varphi^{-1}(0_N) \) of \( \varphi: M \to N \) coincides with the \hyperref[def:categorical_kernels/kernel]{categorical kernel}.

    \thmitem{thm:module_categorical_kernels/cokernel} The module-theoretic cokernel \( \co\ker \varphi = N / \img f \) of \( \varphi: M \to N \) coincides with the \hyperref[def:categorical_kernels/cokernel]{categorical cokernel}.

    \thmitem{thm:module_categorical_kernels/image} The set-theoretic image \( \varphi[M] \) of \( \varphi: M \to N \) coincides with the kernel of the cokernel of \( \varphi \).

    As discussed in \cref{rem:categorical_image_and_coimage}, the image of \( \varphi \) can be defined categorically as
    \begin{equation*}
      \ker \co\ker \varphi \coloneqq \varphi,
    \end{equation*}
    without referring to elements of \( M \) or \( N \).

    \thmitem{thm:module_categorical_kernels/coimage} The set-theoretic coimage \( M / \ker \varphi \) of \( \varphi: M \to N \) coincides with the cokernel of the kernel of \( \varphi \).

    The coimage is briefly discussed in \cref{rem:categorical_image_and_coimage} in the context of \fullref{thm:function_factorization} and related theorems like \fullref{thm:homomorphism_theorem_for_modules}.
  \end{thmenum}
\end{proposition}
\begin{comments}
  \item The statements for cokernels (and hence for images and coimages) do not hold more generally --- see \cref{rem:categorical_image_and_coimage}.
\end{comments}
\begin{proof}
  \SubProofOf{thm:module_categorical_kernels/kernel} Follows from \cref{thm:equational_theory_zero_morphisms/kernel}.

  \SubProofOf{thm:module_categorical_kernels/cokernel} \Cref{thm:equational_theory_zero_morphisms/cokernel} implies that the categorical cokernel of \( \varphi: M \to N \) is the \hyperref[def:fol_congruence]{congruence} \( {\cong} \) \hyperref[def:fol_generated_congruence]{generated} by the relation
  \begin{equation*}
    b \sim 0_N \T{if and only if} b \in \varphi[M],
  \end{equation*}
  endowed with the canonical projection \( \pi: N \to N / {\cong} \).

  Since the generated relation is compatible with sums, this definition can be rewritten slightly by setting \( b = c - d \), so that
  \begin{equation*}
    c \sim d \T{if and only if} c - d \in \varphi[M].
  \end{equation*}

  In fact, the latter relation is itself a congruence because it is compatible with scalar multiplication. The quotient \( N / {\cong} \) thus becomes \( N / \varphi[M] \), which coincides with the module-theoretic cokernel of \( \varphi \) defined in \cref{def:module_cokernel}.

  \SubProofOf{thm:module_categorical_kernels/image} \Cref{thm:module_categorical_kernels/cokernel} implies that the cokernel of \( \varphi: M \to N \) is \( N / \varphi[M] \), endowed with the projection \( \pi: N \to N / \varphi[M] \).

  \Cref{thm:module_categorical_kernels/kernel} implies that the kernel of \( \pi \) is \( \varphi[M] \).

  \SubProofOf{thm:module_categorical_kernels/coimage} \Cref{thm:module_categorical_kernels/kernel} implies that the kernel of \( \varphi: M \to N \) is \( \varphi^{-1}(0_N) \), endowed with the inclusion \( \iota: \varphi^{-1}(0_M) \to M \).

  \Cref{thm:module_categorical_kernels/cokernel} implies that the cokernel of \( \iota \) is \( M / \varphi^{-1}(0_M) \).
\end{proof}

\begin{proposition}\label{thm:module_epimorphisms_are_surjective}
  Every \hyperref[def:fol_epimorphism]{epimorphism} in \hyperref[def:group/category]{\( \cat{R\mathdash Mod} \)} is \hyperref[def:function_invertibility/surjective]{surjective}.
\end{proposition}
\begin{proof}
  Let \( \varphi: M \to N \) be an \( R \)-module epimorphism. Consider the \hyperref[thm:module_categorical_kernels/cokernel]{cokernel} projection \( \pi: N \to N / \img \varphi \) and the parallel \hyperref[def:zero_morphism]{zero morphism} \( 0: N \to N / \img \varphi \). Clearly
  \begin{equation*}
    \pi \bincirc \varphi = 0 \bincirc \varphi.
  \end{equation*}

  Since \( \varphi \) is an epimorphism, it follows that \( \pi = 0 \). Thus, the image \( N / \img \varphi \) of \( \pi \) is the zero module.

  Then, by definition, every element of \( N \) belongs to \( \img \varphi \). This is a restatement of the surjectivity of \( \varphi \).
\end{proof}

\paragraph{Direct sums}

\begin{proposition}\label{thm:internal_module_direct_sum_correctness}\mimprovised
  Suppose that \( N \) and \( K \) are \hyperref[def:module/subobject]{submodules} of \( M \) that intersect at \( \set{ 0_M } \) and whose pointwise sum \( N + K \) is \( M \). Then the \hyperref[def:external_semimodule_direct_sum]{external direct sum} \( N \oplus K \) is isomorphic to \( M \).
\end{proposition}
\begin{proof}
  Our desired isomorphism is the map
  \begin{equation*}
    \begin{aligned}
      &\Phi: N \oplus K \to M, \\
      &\Phi(n, k) \coloneqq n + k.
    \end{aligned}
  \end{equation*}

  \SubProof{Proof that \( \Phi \) is a homomorphism} We have
  \begin{align*}
    \Phi((n_1, k_1) + (n_2, k_2))
    &=
    \Phi(n_1 + n_2, k_1 + k_2)
    = \\ &=
    (n_1 + n_2) + (k_1 + k_2)
    \reloset {\eqref{def:binary_operation/commutative}} = \\ &=
    (n_1 + k_1) + (n_2 + k_2)
    = \\ &=
    \Phi(n_1, k_1) + \Phi(n_2, k_2)
  \end{align*}
  and
  \begin{equation*}
    \Phi(r(n, k))
    =
    \Phi(rn, rk)
    =
    rn + rk
    \reloset {\eqref{eq:def:semimodule/operation/scalar_addition_distributivity}} =
    r(n + k)
    =
    r \cdot \Phi(n, k).
  \end{equation*}

  \SubProofOf[def:function_invertibility/injective/equality]{injectivity} Suppose that \( \Phi(n_1, k_1) = \Phi(n_2, k_2) \). Then
  \begin{equation*}
    (n_1 - n_2) + (k_1 - k_2)
    =
    \Phi(n_1 - n_2, k_1 - k_2)
    =
    \Phi(n_1, k_1) - \Phi(n_2, k_2)
    =
    0_M,
  \end{equation*}
  thus
  \begin{equation*}
    n_1 - n_2 = -(k_1 - k_2).
  \end{equation*}

  The left side belongs to \( N \) and the right side belongs to \( K \), so their common value belongs to \( N \cap K \). But the only value in \( N \cap K \) is \( 0_M \).

  It follows that \( n_1 = n_2 \) and \( k_1 = k_2 \).

  \SubProofOf[def:function_invertibility/surjective/existence]{surjectivity} For any \( x \) in \( M \), since \( M = N + K \), there exist elements \( n \) of \( N \) and \( k \) of \( K \) such that \( g = n + k = \Phi(n, k) \).
\end{proof}

\begin{definition}\label{def:internal_module_direct_sum}\mimprovised
  In the setting of \cref{thm:internal_module_direct_sum_correctness}, we say that \( M \) is an \term[ru=прямая внутренная сумма (\cite[30]{Кострикин2020АлгебраЧасть2}), en=internal direct sum (\cite[60]{Knapp2016BasicAlgebra})]{internal direct sum} of \( N \) and \( K \).

  This generalizes to any finite subfamily of submodules \( N_1, \ldots, N_n \) such that \( N_i \cap N_j = \set{ 0_M } \) for \( i \neq j \) and \( M = N_1 + \ldots + N_n \).
\end{definition}
\begin{comments}
  \item We generalize internal direct sums of two vector subspaces from \cite[60]{Knapp2016BasicAlgebra}.
\end{comments}

\paragraph{Linear dependence}

\begin{definition}\label{def:linear_dependence}
  Let \( M \) be an \( R \)-\hyperref[def:module]{module} and fix a subset \( E \subseteq M \). We say that the elements of \( E \) are \term[bg=линейно независими (вектори) (\cite[94]{Обрешков1962ВисшаАлгебра}), ru=линейно независимые (векторы) (\cite[\S 3.4]{Тыртышников2007ЛинейнаяАлгебра}), en=linearly independent (indexed set) (\cite[def. VI.1.1]{Aluffi2009Algebra})]{linearly independent} if any of the following conditions hold:

  \begin{thmenum}
    \thmitem{def:linear_dependence/direct}\mcite[307]{Aluffi2009Algebra} A \hyperref[def:linear_combination]{linear combination} in \( E \) sums to zero if and only if it is \hyperref[def:free_semimodule]{trivial}.

    \thmitem{def:linear_dependence/evaluation}\mcite[def. VI.1.1]{Aluffi2009Algebra} The \hyperref[thm:free_semimodule_universal_property]{linear combination evaluation map}
    \begin{equation*}
      \begin{aligned}
        &\Phi_E: R^{\oplus E} \to M, \\
        &\Phi_E( \seq{ t_e }_{e \in E} ) \coloneqq \sum_{e \in E} t_e \cdot e
      \end{aligned}
    \end{equation*}
    is \hyperref[def:function_invertibility/injective]{injective}.
  \end{thmenum}

  Unsurprisingly, if the elements of \( E \) are not linearly independent, we say that they are \term{linearly dependent}.
\end{definition}
\begin{comments}
  \item Compare this concept to algebraic dependence defined in \cref{def:algebraic_dependence}.
\end{comments}
\begin{defproof}
  \ImplicationSubProof{def:linear_dependence/direct}{def:linear_dependence/evaluation} Suppose that only the trivial linear combination of \( E \) sums to zero. If \( \sum_{e \in E} t_e \cdot e = \sum_{e \in E} r_e \cdot e \), then
  \begin{equation*}
    \sum_{e \in E} (t_e - r_e) \cdot e
    \reloset {\eqref{eq:def:semimodule/operation/vector_addition_distributivity}} =
    \sum_{e \in E} t_e \cdot e - \sum_{e \in E} r_e \cdot e
    =
    0_M,
  \end{equation*}
  implying that \( t_e = r_e \) for every \( e \) in \( E \). Hence, \( \Phi_E \) is injective.

  \ImplicationSubProof{def:linear_dependence/evaluation}{def:linear_dependence/direct} Trivial.
\end{defproof}

\begin{proposition}\label{thm:def:linear_dependence}
  \hyperref[def:linear_dependence]{Linear (in)dependence} in an \( R \)-module \( M \) has the following basic properties:
  \begin{thmenum}
    \thmitem{thm:def:linear_dependence/empty} The empty set is linearly independent.

    \thmitem{thm:def:linear_dependence/zero} The zero vector \( 0_M \) is by itself linearly dependent.

    \thmitem{thm:def:linear_dependence/monotonicity} If \( E \) is a linearly \emph{dependent} set and \( E \subseteq F \), then \( F \) is also a linearly dependent set.

    \thmitem{thm:def:linear_dependence/antimonotonicity} If \( F \) is a linearly \emph{independent} set and \( E \subseteq F \), then \( E \) is also a linearly independent set.

    \thmitem{thm:def:linear_dependence/dependent_combination} The set \( E \cup \set{ x } \) is linearly dependent if and only if \( x \in \lin E \).

    \thmitem{thm:def:linear_dependence/span_dependent} For any set of vectors \( E \) and any vector \( x \) from their \hyperref[def:linear_span]{linear span}, the set \( E \cup \set{ x } \) is a linearly dependent.

    A related result is \cref{thm:def:vector_space/span_independent}.
  \end{thmenum}
\end{proposition}
\begin{proof}
  \SubProofOf{thm:def:linear_dependence/empty} The empty set cannot be linearly dependent because there are no possible linear combinations of its elements.

  \SubProofOf{thm:def:linear_dependence/zero} Trivial.

  \SubProofOf{thm:def:linear_dependence/monotonicity} Trivial.

  \SubProofOf{thm:def:linear_dependence/antimonotonicity} Trivial.

  \SubProofOf{thm:def:linear_dependence/span_dependent} For any \( x \) in the span of \( E \), by \cref{thm:generated_submodule}, there exists a linear combination of members of \( E \) such that
  \begin{equation*}
    x = \sum_{k=1}^n t_k \cdot e_k.
  \end{equation*}

  If \( x = 0_M \), then \( E \cup \set{ x } \) is linearly dependent by \cref{thm:def:linear_dependence/zero} and \cref{thm:def:linear_dependence/monotonicity}.

  If \( x \neq 0_M \), then \( E \cup \set{ x } \) is linearly dependent because \( x \) is a nontrivial linear combination of other vectors of \( E \).
\end{proof}

\begin{example}\label{ex:def:linear_dependence}
  We list several (counter)examples for \hyperref[def:linear_dependence]{linear (in)dependence}:
  \begin{thmenum}
    \thmitem{ex:def:linear_dependence/irrational_numbers} Like all concepts related to \hyperref[def:linear_combination]{linear combinations}, linear dependence may behave differently depending on the underlying ring. For example, every irrational number is a linear combination of itself, but it is not a linear combination of rational numbers.

    \thmitem{ex:def:linear_dependence/unit_matrix_columns} We have
    \begin{equation*}
      2 \cdot
      \begin{pmatrix}
        1 \\ 3
      \end{pmatrix}
      =
      \begin{pmatrix}
        2 \\ 0
      \end{pmatrix}
      +
      3 \cdot
      \begin{pmatrix}
        0 \\ 2
      \end{pmatrix}
    \end{equation*}

    Thus, the three column vectors are linearly dependent.

    The above equality holds for matrices over any ring, however only if \( 2 \) is invertible we can express the first vector via the others:
    \begin{equation*}
      \begin{pmatrix}
        1 \\ 3
      \end{pmatrix}
      =
      \frac 1 2 \cdot
      \begin{pmatrix}
        2 \\ 0
      \end{pmatrix}
      +
      \frac 3 2 \cdot
      \begin{pmatrix}
        0 \\ 2
      \end{pmatrix}
    \end{equation*}
  \end{thmenum}
\end{example}

\paragraph{Hamel bases}

\begin{definition}\label{def:hamel_basis}
  Let \( M \) be a left \hyperref[def:module]{\( R \)-module} and fix a subset \( E \subseteq M \). We say that \( E \) is a \term{Hamel basis} or simply \term[bg=базис (\cite[100]{Обрешков1962ВисшаАлгебра}), ru=базис (\cite[\S 3.7]{Тыртышников2007ЛинейнаяАлгебра}), en=Hamel basis (\cite[example B-2.12]{Rotman2015AdvancedModernAlgebraPart1})]{basis} of \( M \) if any of the following equivalent conditions hold:

  \begin{thmenum}
    \thmitem{def:hamel_basis/independent}\mcite[307]{Aluffi2009Algebra} It is a \hyperref[thm:generated_submodule]{spanning set} of \hyperref[def:linear_dependence]{linearly independent} elements.

    \thmitem{def:hamel_basis/evaluation}\mcite[lemma VI.1.5]{Aluffi2009Algebra} The \hyperref[thm:free_semimodule_universal_property]{free semimodule evaluation map} is \hyperref[def:function_invertibility/bijective]{bijective}.

    \thmitem{def:hamel_basis/free}\mimprovised \( M \) is \hyperref[def:linear_function]{linearly isomorphic} to the \hyperref[def:free_semimodule]{free module} \( R^{\oplus E} \).
  \end{thmenum}
\end{definition}
\begin{comments}
  \item In relation to \cref{def:hamel_basis/free}, we may refer to \( M \) itself is a free module.
  \item By \fullref{thm:vector_space_basis_existence}, A Hamel basis always exists for \hyperref[def:vector_space]{vector spaces}, but \cref{ex:def:hamel_basis/module_without_basis} shows that a more general module may fail to have a basis.

  Even in the latter case, \fullref{thm:fundamental_theorem_of_finitely_generated_abelian_groups} demonstrates how a finitely-generated \hyperref[def:abelian_group]{abelian group} \( G \) has a decomposition into \hyperref[def:cyclic_group]{cyclic groups}. This is generally not a \( \BbbZ \)-basis of \( G \) unless all groups are infinite.
\end{comments}
\begin{defproof}
  \ImplicationSubProof{def:hamel_basis/independent}{def:hamel_basis/evaluation} Suppose that \( E \) is a spanning set of linearly independent elements.

  \Cref{def:linear_dependence/evaluation} is satisfied, hence the evaluation map is injective.

  Furthermore, \( E \) spans \( M \); hence, the evaluation map is also surjective.

  \ImplicationSubProof{def:hamel_basis/evaluation}{def:hamel_basis/free} If the evaluation map is bijective, it is itself a linear isomorphism.

  \ImplicationSubProof{def:hamel_basis/free}{def:hamel_basis/independent} Suppose that \( \Psi: R^{\oplus E} \to M \) is a linear isomorphism. Then \( \Psi \) is surjective, and hence a spanning set of \( M \). Furthermore, it satisfies \cref{def:linear_dependence/evaluation}, meaning that \( E \) is a linearly independent set.
\end{defproof}

\begin{remark}\label{rem:free_module}
  After discussing \hyperref[con:free_construction]{free constructions} in generality, and directly after showing that an \hyperref[def:module]{\( R \)-module} with a \hyperref[def:hamel_basis]{Hamel basis} \( \set{ e_1, \ldots, e_n } \)  is isomorphic to the \hyperref[def:free_semimodule]{free semimodule} \( R^n \), Nathan Jacobson writes in \cite[171]{Jacobson1985BasicAlgebraI} the following:
  \begin{displayquote}
    We have therefore shown that the existence of a base of \( n \) elements for a module \( M \) implies that \( M \cong R^{(n)} \). In this case we shall say that \( M \) is a \textit{free \( R \)-module of rank \( n \)}.
  \end{displayquote}

  We generalize this to \hyperref[con:transfinitum]{transfinite} bases --- we will refer to any module with a basis as \enquote{free} since choosing a Hamel basis \( E \) for \( M \) gives us an explicit isomorphism from \( M \) to \( R^{\oplus E} \). This is in line with the practice of identifying isomorphic objects discussed in \cref{rem:free_construction_univalence}.

  We will define the rank of a free module in \cref{def:module_rank}, after proving some coherence results.
\end{remark}
\begin{comments}
  \item The above convention is also used by other authors. A \enquote{free \( R \)-module} with basis \( E \) is defined as a module isomorphic to \( R^{\oplus A} \) in
  \cite[def. II.1.10]{Bourbaki1998Algebra1to3},
  \cite[135]{Lang2002Algebra},
  \cite[329]{Rotman2015AdvancedModernAlgebraPart1},
  \cite[169]{Aluffi2009Algebra},
  \cite[405]{Knapp2016BasicAlgebra} and
  \cite[47]{Шафаревич2019ОсновныеПонятияАлгебры}.
\end{comments}

\begin{example}\label{ex:def:hamel_basis}
  We list examples of \hyperref[def:hamel_basis]{Hamel bases}:
  \begin{thmenum}
    \thmitem{ex:def:hamel_basis/euclidean_space} The canonical example of a basis is the following set of column vectors of the \hyperref[eq:thm:matrix_algebra/matrix_multiplication/identity]{identity matrix} in the \hyperref[def:euclidean_space]{Euclidean space} \( \BbbR^n \), that is, the vectors
    \begin{equation*}
      \begin{pmatrix}
        1 \\ 0 \\ \ldots \\ 0 \\ 0
      \end{pmatrix},
      \begin{pmatrix}
        0 \\ 1 \\ \ldots \\ 0 \\ 0
      \end{pmatrix},
      \cdots,
      \begin{pmatrix}
        0 \\ 0 \\ \ldots \\ 0 \\ 1
      \end{pmatrix}.
    \end{equation*}

    \thmitem{ex:def:hamel_basis/polynomial_algebra} In \cref{def:polynomial_algebra}, we have defined a polynomial algebra as a \hyperref[def:semigroup_algebra]{monoid algebra} over the set of all monomials; a monoid algebra is in turn defined as a \hyperref[def:free_semimodule]{free \( R \)-module} over the monomials with a convolution product.

    Hence, for a nontrivial commutative ring \( R \) and for indeterminates \( \mscrX \), the set
    \begin{equation*}
      \set*{ X_1 \cdots X_n \given X_1, \ldots, X_n \in \mscrX }
    \end{equation*}
    of all monomials is a \hyperref[def:hamel_basis]{basis} for \hyperref[def:polynomial_algebra]{polynomial ring} \( R[\mscrX] \).

    \thmitem{ex:def:hamel_basis/module_without_basis} As a \hyperref[thm:abelian_group_is_module]{\( \BbbZ \)-module}, the additive group of \( \BbbQ \) has no basis.

    Indeed, given any two rational numbers \( a / b \) and \( c / d \), we have
    \begin{equation*}
      cb \cdot \frac a b + (-da) \cdot \frac c d = 0.
    \end{equation*}

    Therefore, only singleton sets of rational numbers are linearly independent with respect to \( \BbbZ \). But no single integer generates \( \BbbQ \).
  \end{thmenum}
\end{example}

\begin{proposition}\label{thm:def:hamel_basis}
  \hyperref[def:hamel_basis]{Hamel bases} of the \( R \)-module \( M \) have the following basic properties:
  \begin{thmenum}
    \thmitem{thm:def:hamel_basis/span} Every \hyperref[def:linear_dependence]{linearly independent} subset of \( M \) is a basis for its \hyperref[def:semimodule/subobject]{linear span}.

    \thmitem{thm:def:hamel_basis/maximal_independent} Every basis of \( M \) is maximal among \hyperref[thm:def:linear_dependence]{linearly independent} sets.

    The converse holds for vector spaces --- see \cref{thm:def:vector_space/minimal_spanning}.

    \thmitem{thm:def:hamel_basis/minimal_spanning} Every basis of \( M \) is minimal among \hyperref[def:linear_span]{spanning sets}.

    The converse holds for vector spaces --- see \cref{thm:def:vector_space/minimal_spanning}.
  \end{thmenum}
\end{proposition}
\begin{proof}
  \SubProofOf{thm:def:hamel_basis/span} Trivial.

  \SubProofOf{thm:def:hamel_basis/maximal_independent} Follows from \cref{thm:def:linear_dependence/span_dependent}.

  \SubProofOf{thm:def:hamel_basis/minimal_spanning} Let \( E \) be a basis of \( M \) and suppose that \( F \subsetneq E \) is also a spanning set of \( M \).

  Then there exists some vector \( x \in E \setminus F \). Since both sets are spanning, \( \lin F = \lin E \). By \cref{thm:def:linear_dependence/span_dependent}, \( F \cup \set{ x } \) is linearly dependent, and by \cref{thm:def:linear_dependence/antimonotonicity}, \( E \) is linearly dependent. This contradicts the assumption that \( E \) has a basis.

  Therefore, no proper subset of \( E \) can span \( M \).
\end{proof}

\begin{definition}\label{def:ordered_hamel_basis}\mcite[79]{FriedbergInselSpence2018LinearAlgebra}
  We refer to any \hyperref[def:set_enumeration]{enumeration} of a \hyperref[def:hamel_basis]{Hamel basis} as an \term{ordered basis}.
\end{definition}

\begin{proposition}\label{thm:basis_of_direct_sum}
  Consider the family of \( R \)-modules \( \seq{ M_k }_{k \in \mscrK} \). Suppose that \( \seq{ E_k }_{k \in \mscrK} \) are bases of the corresponding modules. Then the union
  \begin{equation*}
    E \coloneqq \bigcup_{k \in \mscrK} \iota_k[E_k]
  \end{equation*}
  is a basis for the \hyperref[def:external_semimodule_direct_sum]{direct sum}
  \begin{equation*}
    M \coloneqq \bigoplus_{k \in \mscrK} M_k.
  \end{equation*}
\end{proposition}
\begin{comments}
  \item For the sake of preciseness, here we distinguish between \( M_k \) and its embedding \( \iota_k[E_k] \) into \( M \).
\end{comments}
\begin{proof}
  As subspaces of \( M \), \( \iota_k[M_k] \) and \( \iota_m[M_m] \) are disjoint for \( k \neq m \). Thus, every vector from \( M \) is a unique sum of vectors from the subspaces.

  Since, for every \( k \) in \( \mscrK \), every vector from \( \iota_k[M_k] \) can further be uniquely represented as a linear combination of vectors from \( \iota_k[E_k] \), we conclude that every vector from the direct sum is a linear combination of vectors from \( E \).
\end{proof}

\begin{definition}\label{def:basis_decomposition}\mimprovised
  If \( E \) is a \hyperref[def:hamel_basis]{Hamel basis} of the \( R \)-\hyperref[def:module]{module} \( M \), the inverse of the linear isomorphism from \cref{def:hamel_basis/evaluation} is
  \begin{equation*}
    \begin{aligned}
      &\pi_E: M \to R^{\oplus E} \\
      &\pi_E\parens*{ \sum_{e \in E} t_e \cdot e } \coloneqq \seq{ t_e }_{e \in E}.
    \end{aligned}
  \end{equation*}

  We denote the \( e \)-th component of this function by \( \pi_e: M \to R \). It is a linear map, which we call the \term[en=coordinate projection (\cite[\S 1.34]{Schechter1997AnalysisHandbook})]{coordinate projection} onto \( e \).

  For every vector \( x \) in \( M \), we thus have
  \begin{equation*}
    x = \sum_{e \in E} \pi_e(x) \cdot e.
  \end{equation*}

  This linear combination is unique, and we call it the \term[ru=разложение по базису (\cite[sec. 12.1]{Тыртышников2007ЛинейнаяАлгебра})]{basis decomposition} of \( x \) along \( E \).
\end{definition}
\begin{comments}
  \item As in \cref{def:linear_combination}, we sometimes take only the basis vectors with nonzero coefficients as \enquote{the} decomposition.

  \item We have defined projection maps of \hyperref[def:cartesian_product]{Cartesian products} in \cref{def:cartesian_product_projection} as functions from \( R^{\oplus E} \) to \( R \). Here the coordinate projections generalize this by taking \( M \) as its domain.
\end{comments}

\begin{proposition}\label{thm:basis_projection_orthonormal}
  Let \( E \) be a \hyperref[def:hamel_basis]{Hamel basis} of the \hyperref[def:free_semimodule]{free} \( R \)-\hyperref[def:module]{module} \( M \). Then the \hyperref[def:basis_decomposition]{projection functionals} satisfy
  \begin{equation*}
    \pi_e(f) = \delta_{e,f} = \begin{cases}
      1_R, &e = f, \\
      0_R, &\T{otherwise.}
    \end{cases}
  \end{equation*}
\end{proposition}
\begin{proof}
  We have
  \begin{equation*}
    f = \sum_{e \in B} \pi_e(f) \cdot e.
  \end{equation*}

  Thus,
  \begin{equation*}
    0_M = f - \sum_{e \in B} \pi_e(f) \cdot e = (1 - \pi_f(f)) \cdot f - \sum_{e \in E \setminus \set{ f }} \pi_e(f).
  \end{equation*}

  Since the vectors in \( E \) are linearly independent, \( \pi_f(f) = 1 \) and \( \pi_e(f) = 0_M \) for \( e \in E \).
\end{proof}

\begin{definition}\label{def:coordinate_space}\mimprovised
  For a \hyperref[def:ring/commutative]{commutative ring} \( R \), consider the \( n \)-fold \hyperref[def:cartesian_product]{Cartesian product} \( R^n \), the \hyperref[def:free_semimodule]{free module} consisting of \hyperref[def:ordered_tuple]{ordered tuples} of length \( n \). We call \( R^n \) a \term[ru=координатное пространство (\cite[443]{ИльинСадовничийСендов1985АнализТом1}), en=coordinate space (\cite[4]{Halmos1974VectorSpaces})]{coordinate space}.

  \Cref{thm:basis_of_direct_sum} shows that the vectors \( e_1, \ldots, e_n \), in which the \( i \)-th component of the tuple \( e_i \) is \( 1 \) and all others are \( 0_M \), is an \hyperref[def:ordered_hamel_basis]{ordered} \hyperref[def:hamel_basis]{Hamel basis}, which we call the \term{standard basis} of \( R^n \).
\end{definition}
\begin{comments}
  \item Since the members of \( R^n \) are \hyperref[def:ordered_tuple]{ordered tuples}, and we often regard them as \hyperref[def:array/column_vector]{column vectors}. See \cref{rem:vector} for a more detailed discussion of terminology for vectors.

  \item We generalize the real and complex coordinate spaces discussed by \textcite[4; 5]{Halmos1974VectorSpaces}.
\end{comments}

\paragraph{Torsion}

\begin{definition}\label{def:module_torsion}\mcite[def. VI.4.1]{Aluffi2009Algebra}
  We say that the vector \( x \) in an \hyperref[def:module]{\( R \)-module} is a \term[ru=элемент кручения (\cite[57]{Шафаревич2019ОсновныеПонятияАлгебры})]{torsion element} if there exists some nonzero scalar \( t \) such that \( tx \) is the zero vector. A module without torsion elements is called \term{torsion-free}.
\end{definition}

\begin{example}\label{ex:def:module_torsion}
  We list examples of \hyperref[def:module_torsion]{torsion elements}:
  \begin{thmenum}
    \thmitem{ex:def:module_torsion/scalars} When regarding a ring as a (bi)module over itself, every \hyperref[def:semiring_divisibility/zero]{proper zero divisor} contains is a torsion element. For example, in \( \BbbZ_6 \),
    \begin{equation*}
      2 \cdot 3 = 3 \cdot 2 = 0 \pmod 6,
    \end{equation*}
    hence both \( 2 \) and \( 3 \) are torsion elements.

    \thmitem{ex:def:module_torsion/matrices} The matrix
    \begin{equation*}
      \begin{pmatrix}
        2 & 0 \\
        0 & 4
      \end{pmatrix}
    \end{equation*}
    is a torsion element of the \hyperref[thm:matrix_algebra]{matrix algebra} \( \BbbZ_6^{2 \times 2} \) because
    \begin{equation*}
      3 \cdot 2 = 3 \cdot 4 = 0 \pmod 6.
    \end{equation*}
  \end{thmenum}
\end{example}

\begin{proposition}\label{thm:basis_implies_torsion_free}
  Every \hyperref[def:free_semimodule]{free module} is \hyperref[def:module_torsion]{torsion-free}.
\end{proposition}
\begin{proof}
  Let \( E \) be a basis of the \( R \)-module \( M \).

  Suppose that \( tx \) for some nonzero scalar \( t \) and nonzero vector \( x \). Then
  \begin{equation*}
    0_R = tx = \sum_{e \in E} (t \pi_e(x)) \cdot e
  \end{equation*}
  is a nontrivial linear combination that sums to zero, contradicting the assumption that \( E \) is a basis.

  Therefore, \( x \) cannot be a torsion element.
\end{proof}
